\documentclass[11pt]{article}
\usepackage[margin=0.8in]{geometry}
\usepackage[T1]{fontenc}
\usepackage{lmodern,amsmath,amssymb,amsthm,booktabs,array}
\usepackage[colorlinks=true,linkcolor=blue,urlcolor=blue,citecolor=blue]{hyperref}
\hypersetup{pdftitle={Four sequence proofs: congruences, growth, and Euler products},pdfauthor={Blake Stone},pdfsubject={OEIS conjectures with Lean verification}}
\setlength{\parindent}{0pt}
\setlength{\parskip}{0.4em}
\newtheorem{theorem}{Theorem}[section]
\newtheorem{lemma}[theorem]{Lemma}
\newcommand{\N}{\mathbb N}
\newcommand{\Z}{\mathbb Z}
\newcommand{\R}{\mathbb R}
\begin{document}
\begin{center}
{\LARGE\bfseries Four sequence proofs:\\[3pt]
congruences, growth, and Euler products}\par
\vspace{0.5em}
{\large Blake Stone}\par
Proof developed with OpenAI Astra 6.2\par
November 26, 2026
\end{center}

\textbf{Abstract.}
We prove five assertions recorded as conjectural or tentative in four OEIS
entries. We count the ordered solutions of a two-cube congruence modulo
powers of three, classify the residue of a weighted binomial sum by its
index's base-five digits, establish a golden-ratio limit for a complementary
recurrence with a negative correction, and identify both the columns and
row sums of a recursively defined triangle. The accompanying five Lean
modules have fresh package builds and independent compilations. The
recurrence theorem derives its growth hypotheses from the defining
conditions; it does not separately formalize existence or uniqueness of
the complementary pair. Background algebra and earlier related arguments
are credited, and no exhaustive historical-priority claim is made.

\section{A087412: a cubic congruence modulo powers of three}

Let
\[
 a(N)=\#\{(x,y)\in\{0,\ldots,N-1\}^2:
                 x^3+y^3\equiv1\pmod N\}.
\]
The pairs are ordered. Andrew Howroyd's July 17, 2018 conjecture in
A087412~\cite{A087412} concerns the prime $3$.

\begin{theorem}\label{thm:cubes}
For every integer $e\ge2$, $a(3^e)=2\cdot3^e$.
\end{theorem}

\begin{proof}
Write $e=k+2$ and $M=3^k$. Each residue modulo $9M$ is uniquely $r+3u$,
where $r\in\{0,1,2\}$ and $0\le u<3M$. Since
$(r+3u)^3\equiv r^3\pmod9$, the possible cubic residues are $0,1,8$.
Two of these sum to $1$ modulo $9$ only in the ordered patterns $(0,1)$
and $(1,0)$. Thus every solution has one of the disjoint forms
\[
 (x,y)=(3u,1+3v)\quad\hbox{or}\quad(1+3u,3v),
 \qquad 0\le u,v<3M.
\]
Swapping the coordinates exchanges the two cases.

Consider the polynomial $F(t)=t+3t^2+3t^3$. Its difference factors as
\[
 F(s)-F(t)=(s-t)\bigl[1+3(s+t+s^2+st+t^2)\bigr].
\]
For $k>0$, the bracket is a unit modulo $3^k$, because it is congruent to
$1$ modulo $3$. Consequently $F$ is injective, hence bijective, on the
finite ring $\Z/M\Z$. For $k=0$ the ring has one element, so bijectivity
also holds.

In the first residue pattern, expansion gives
\[
 (3u)^3+(1+3v)^3-1=9\bigl(3u^3+F(v)\bigr).
\]
The original congruence is therefore equivalent to
$F(v)\equiv-3u^3\pmod M$. For each of the $3M$ possible values of $u$,
bijectivity gives exactly one solution $v_0$ modulo $M$. Its representatives
in $0\le v<3M$ are precisely $v_0,v_0+M,v_0+2M$. This gives $9M$ pairs
in the first case and another $9M$ in the second. The total is
$18M=2\cdot3^e$.
\end{proof}

The boundary condition matters: $a(3)=3$, while the first covered values
are $a(9)=18$, $a(27)=54$, and $a(81)=162$. In Lean,
\texttt{A087412.original\_count} identifies the indicator sum with the
cardinality of the literal ordered-pair solution set, and
\texttt{A087412.original\_conjecture} states Theorem~\ref{thm:cubes}
directly in terms of that cardinality. The unique residue decomposition,
permutation argument, and three lifts per root are all formalized.
The unit argument is implemented using nilpotence of $3$ in
$\Z/3^k\Z$; it covers the one-element ring as well.

\section{A217280: the zero-or-two base-five classification}

Paul D. Hanna's A217280~\cite{A217280}, introduced September 29, 2012,
is the weighted binomial sum
\[
 a(n)=\sum_{j=0}^{n}\binom nj^3 2^j3^{n-j},\qquad n\ge0.
\]

\begin{theorem}\label{thm:digits}
The remainder of $a(n)$ modulo $5$ is $1$ if every digit in the base-five
representation of $n$ is $0$ or $2$, and is $0$ otherwise.
\end{theorem}

\begin{proof}
Write $n=5q+r$, where $0\le r<5$. Decompose the summation index as
$j=5h+s$, where $0\le s<5$. Lucas's theorem gives
\[
 \binom{5q+r}{5h+s}\equiv\binom qh\binom rs\pmod5.
\]
If $s>r$, the factor $\binom rs$ is zero. If $s\le r$ and $h\le q$,
there is no borrow:
\[
 5q+r-(5h+s)=5(q-h)+(r-s).
\]
Fermat's congruences $2^5\equiv2$ and $3^5\equiv3\pmod5$ therefore yield
\begin{align*}
 &\binom{5q+r}{5h+s}^{\!3}2^{5h+s}3^{5q+r-5h-s}\\
 &\hspace{1cm}\equiv
 \left(\binom qh^3 2^h3^{q-h}\right)
 \left(\binom rs^3 2^s3^{r-s}\right)\pmod5.
\end{align*}
To justify factoring the finite sum, pad it with zero summands for
$j>5q+r$ until its range is $0\le j<5(q+1)$. This convention avoids
negative exponents in the padded terms; their binomial coefficients
are zero. The new range is the rectangle
$0\le h\le q$, $0\le s<5$. The terms with $s>r$ vanish modulo $5$,
so the sum factors to give
\begin{equation}\label{eq:lucas}
 a(5q+r)\equiv a(q)a(r)\pmod5.
\end{equation}

The single-digit values and remainders are
\[
\begin{array}{c|rrrrr}
 r&0&1&2&3&4\\\hline
 a(r)&1&5&61&845&12865\\
 a(r)\bmod5&1&0&1&0&0.
\end{array}
\]
Apply~\eqref{eq:lucas} repeatedly, removing the last base-five digit at
each step. The remainder is the product of the corresponding digit
remainders. This product is $1$ exactly when all digits are $0$ or $2$,
and is $0$ otherwise. For $n=0$, the empty-digit convention gives the
empty product $1$, agreeing with $a(0)=1$.
\end{proof}

For example, $62=(222)_5$ gives remainder $1$, whereas $63=(223)_5$
gives remainder $0$. The formal proof starts from the displayed sum,
proves the padding and rectangular reindexing, and applies mathlib's
Lucas theorem~\cite{Lucas}. Its digit predicate uses
\texttt{Nat.digits 5 n} directly. The public declarations
\texttt{A217280.conjecture}, \texttt{one\_iff}, and
\texttt{divisible\_iff} cover every natural index. The Lucas method and
Fermat congruences are classical; the target is this particular
classification, not a new general digit-congruence principle.

\section{A294413: growth despite a negative complementary correction}

Let $a,b:\N\to\N$ be strictly increasing sequences with disjoint ranges
covering the positive integers. Require
\[
 a_0=1,\quad a_1=3,\quad b_0=2,\quad b_1=4,
 \qquad a_{n+2}=a_{n+1}+a_n-b_{n+1}+6.
\]
The equation is interpreted over the integers. These are the defining
conditions in the title and program of Clark Kimberling's
A294413~\cite{A294413}.

\begin{theorem}\label{thm:ratio}
For every pair satisfying these conditions, with
$\phi=(1+\sqrt5)/2$,
\[
 \lim_{n\to\infty}\frac{a_{n+1}}{a_n}=\phi.
\]
More quantitatively, for every $n\ge0$,
\[
 \left|\frac{a_{n+2}}{a_{n+1}}-\phi\right|
 \le20(n+1)\left(\frac23\right)^n.
\]
\end{theorem}

\begin{proof}
We first establish growth from the defining recurrence. If $t<a_i$ and
$t<b_j$, each integer from $1$ through $t$ occurs among the first $i$
terms of $a$ or the first $j$ terms of $b$. Hence
\begin{equation}\label{eq:threshold}
 t\le i+j.
\end{equation}
It follows that $a_i>i+j+1$ implies $b_j\le a_i$: otherwise apply
\eqref{eq:threshold} at threshold $a_i$ with indices $i+1,j$.

We claim $a_n\ge3n$. The initial cases follow from
$a_0=1$, $a_1=3$, and $a_2=6$. For $n\ge3$, strong induction gives
$a_{n-1}\ge3(n-1)>2(n-1)+1$. The preceding consequence with
$i=j=n-1$ gives $b_{n-1}\le a_{n-1}$. Thus
\[
 a_n=a_{n-1}+a_{n-2}-b_{n-1}+6
 \ge a_{n-2}+6\ge3n.
\]
Next, if $b_n>2n+2$, then $a_{n+1}\ge3n+3>2n+2$, and
\eqref{eq:threshold} at indices $n+1,n$ gives the impossible inequality
$2n+2\le2n+1$. We have therefore proved
\begin{equation}\label{eq:complement}
 b_n\le2n+2.
\end{equation}

Subtract a linear part by defining $u_n=a_{n+1}-2(n+1)$. Its initial
values are $u_0=1$, $u_1=2$, and substitution gives
\[
 u_{n+2}=u_{n+1}+u_n+2n+6-b_{n+2}.
\]
The last term is nonnegative by~\eqref{eq:complement}. Since
$(3/2)^2\le3/2+1$, two-step induction proves
\begin{equation}\label{eq:growth}
 u_n\ge(3/2)^n,\qquad a_{n+1}\ge(3/2)^n.
\end{equation}

Now put $A_n=a_{n+1}$ and $e_n=6-b_{n+2}$. The recurrence becomes
$A_{n+2}=A_{n+1}+A_n+e_n$, with $|e_n|\le6(n+1)$.
For the step error $d_n=A_{n+1}-\phi A_n$, the identity
$\phi^2=\phi+1$ gives
\[
 d_{n+1}=(1-\phi)d_n+e_n.
\]
Since $|1-\phi|\le2/3$, setting $K=18+|d_0|$ yields
$|d_n|\le K(n+1)$ by induction: the next bound is at most
$\tfrac23K(n+1)+6(n+1)\le K(n+1)\le K(n+2)$.
Here $d_0=6-3\phi$ and $|6-3\phi|\le2$, so $K\le20$.
Dividing by the positive lower bound~\eqref{eq:growth} proves the
displayed quantitative estimate. Its right side tends to zero. Shifting
the ratio sequence by one index leaves its limit unchanged, which proves
the original conjecture.
\end{proof}

\textbf{Source and prior-work qualifications.}
The source's example complement incorrectly includes $6$, although
$a_2=6$; the title, seeds, program, and $a$-data agree with the conditions
above. The actual complement starts $2,4,5,7,8,9,11,\ldots$.
The formal theorem quantifies over all pairs satisfying the conditions
and assumes neither growth nor convergence. It does not separately
construct the mex sequences or prove existence and uniqueness.

Theorem 2.1 of Kimberling and Moses~\cite{KM} assumes existence of the
ratio limit and a positive disturbance, so it does not prove this case.
An earlier public note~\cite{Note} claims the same case through a general
signed-disturbance theorem. That theorem as stated has a counterexample:
$a_n\equiv1$ satisfies $a_n=a_{n-1}+a_{n-2}-1$, remains positive, and has
$\sum_{n\ge2}\phi^{-n}<\infty$, but its ratios equal $1$.
The missing growth condition is supplied here by~\eqref{eq:growth}.
Related positive-disturbance arguments in that note are prior work. We
claim the exact formal completion and derived growth bound, without
claiming invention of recurrence normalization or exhaustive priority.

\section{A380385: Euler-product columns and row sums}

Mats Granvik's A380385~\cite{A380385} defines a triangle by a recursive
program and gives two tentative identifications. Its explicitly designated
program sets $T(n,1)=1$ for $n=1$ and $T(n,1)=0$ otherwise. For $k\ge2$,
it sets $T(n,k)=0$ when $n<k$, and otherwise
\begin{equation}\label{eq:triangle}
 T(n,k)=\sum_{i=1}^{k-1}T(n-i,k-1)
        -\sum_{i=1}^{n-1}T(n-i,k).
\end{equation}
The first-column exception is essential and agrees with the displayed
data. The shorter title recurrence omits it; applying that title literally
to the first column would give a different array.

\begin{theorem}\label{thm:triangle}
For $1\le k\le n$,
\[
 T(n,k)=[x^{n-k}]\prod_{i=1}^{k-1}(1-x^i).
\]
For every $n\ge1$,
\[
 \sum_{k=1}^{n}T(n,k)=-\operatorname{A010815}(n),
\]
where A010815 is the coefficient sequence of
$\prod_{i\ge1}(1-x^i)$.
\end{theorem}

\begin{proof}
Extend $T(n,k)=0$ for negative $n$. All generating series are formal. Put
$C_k(x)=\sum_{n\ge0}T(n,k)x^n$. The exceptional first column gives
$C_1(x)=x$. Move the second sum in~\eqref{eq:triangle} to the left to obtain
\[
 \sum_{j=1}^{n}T(j,k)=\sum_{i=1}^{k-1}T(n-i,k-1).
\]
The same equation holds when $n<k$, since all relevant terms vanish.
Taking generating series yields
\[
 \frac{C_k(x)}{1-x}=(x+x^2+\cdots+x^{k-1})C_{k-1}(x).
\]
Consequently $C_k(x)=x(1-x^{k-1})C_{k-1}(x)$, and induction gives
\[
 C_k(x)=x^k\prod_{i=1}^{k-1}(1-x^i).
\]
Extracting the coefficient of $x^n$ proves the column formula.

For the row identity, write $P_m(x)=\prod_{i=1}^m(1-x^i)$, with
$P_0(x)=1$. The proved column polynomial satisfies
\[
 C_k(x)=x^kP_{k-1}(x)=P_{k-1}(x)-P_k(x).
\]
Thus $\sum_{k=1}^{N}C_k(x)=1-P_N(x)$. Choose $N=n\ge1$ and extract
the coefficient of $x^n$ to obtain
\[
 \sum_{k=1}^{n}T(n,k)=-[x^n]P_n(x).
\]
Every factor $1-x^i$ with $i>n$ leaves this coefficient unchanged.
It is therefore the degree-$n$ coefficient of the infinite Euler product,
which is precisely A010815~\cite{A010815}.
\end{proof}

The formal proof defines the recursive triangle first. Strong induction
then establishes the column formula, using finite polynomial arithmetic.
The theorem \texttt{A380385.eulerCoefficient\_stable} supplies the exact
connection between finite products and the infinite-product coefficient;
no analytic convergence premise is needed. Euler's product, its known
coefficient sequence, and the related finite-product triangle A231599 are
prior mathematics. The two results concern their asserted connection to
this particular recurrence.

\section{Verification, scope, and source record}

The repository contains four standalone packages, with five Lean modules
in total. Each pins Lean 4.19.0 and mathlib commit
\texttt{c44e0c8ee63ca166450922a373c7409c5d26b00b}. Each package has a fresh
successful build, a separate mathematical/source review, and independent
compilation of the exact final sources. Source hashes connect these checks
to the delivered files. The final theorem axiom reports contain only
\texttt{propext}, \texttt{Classical.choice}, and \texttt{Quot.sound}; there
are no proof placeholders, additional axioms, or \texttt{native\_decide}
calls.

\begin{center}
\begin{tabular}{@{}ll@{}}
\toprule
Package & Main formal declarations\\
\midrule
A087412 & \texttt{original\_conjecture}\\
A217280 & \texttt{conjecture}, \texttt{one\_iff}, \texttt{divisible\_iff}\\
A294413 & \texttt{conjecture}, \texttt{ratio\_error\_bound}\\
A380385 & \texttt{column\_formula}, \texttt{row\_sum}\\
\bottomrule
\end{tabular}
\end{center}

The source snapshots retain the selected assertions as conjectural or
tentative. Exact-identifier and mathematical-description searches, the
cited literature, and public proof collections were checked. The earlier
attempt concerning A294413 is acknowledged explicitly. These are bounded
source-status findings, not proof of absolute historical priority or
evidence of OEIS editorial acceptance. The batch resolves five selected
assertions across four entries, using classical methods with complete
formal verification.

The A294413 analytic module is reused from the repository's A294532
package; its signed-error result explicitly requires the lower growth
bound established here. A separate prior-work addendum in the A294532
directory records the related notes found during this follow-up search.

Repository: \url{https://github.com/blakestone-x/math}. The package
READMEs contain full reproduction instructions, precise source-definition
qualifications, and links to the audit and compilation records.

\begin{thebibliography}{9}
\raggedright
\bibitem{A087412}
OEIS Foundation, A087412, sequence by Yuval Dekel (2003), selected
conjecture by Andrew Howroyd (July 17, 2018).
\url{https://oeis.org/A087412}. Official export revision 30,
September 17, 2019.

\bibitem{A217280}
P. D. Hanna, OEIS A217280 (September 29, 2012).
\url{https://oeis.org/A217280}. Official export revision 25,
July 27, 2022.

\bibitem{Lucas}
G. Ma and mathlib contributors, \emph{Lucas's theorem}, formalization in
\texttt{Mathlib/Data/Nat/Choose/Lucas.lean}, pinned mathlib version 4.19.0.
\url{https://github.com/leanprover-community/mathlib4}.

\bibitem{A294413}
C. Kimberling, OEIS A294413 (October 31, 2017).
\url{https://oeis.org/A294413}. Official export revision 9,
November 1, 2017.

\bibitem{KM}
C. Kimberling and P. J. C. Moses, \emph{Linear Complementary Equations and
Systems}, The Fibonacci Quarterly 57(5) (2019), 96--110.
\url{https://www.fq.math.ca/Papers1/57-5/kimberling-moses.pdf}.

\bibitem{Note}
\emph{Kimberling Note}, publicly uploaded by \texttt{bambasene564};
publication date not established. Retrieved November 26, 2026.
\url{https://www.scribd.com/document/1067832042/kimberling-note}.

\bibitem{A380385}
M. Granvik, OEIS A380385 (January 23, 2025), with coefficient program
dated January 25, 2025. \url{https://oeis.org/A380385}.
Official export revision 20, January 25, 2025.

\bibitem{A010815}
OEIS Foundation, A010815, Euler-product coefficients.
\url{https://oeis.org/A010815}. Definition checked against the official
export, revision 271, May 23, 2026.
\end{thebibliography}
\end{document}
